% v9. A genuinely geometric local high-denominator example for four lines.
% Uses existing bibliography key EM and per:chart-definition.

\section{四直線 core の奇数分母と五次の追加制約}
\label{q9:quintic-local-high-q}

四 puncture neck の剛性は両端の比較から生じる．
三 puncture 曲線だけを用いる四直線 core では，高分母の作用適合 chart を実際に構成できる．
次の基準では，局所の整数条件と正の basin areas が十分であることを，sector 全体で確認する．

\begin{proposition}[四直線 core の十分条件]\label{q9:four-line-criterion}
$L_1,\ldots,L_4\subset\PP^2$ を一般位置の四直線とし，正の weights $w_i$ と
\[
 [\omega]=2\pi\sum_{i=1}^4 w_i[L_i],\qquad W=\sum_iw_i
\]
を固定する．各直線に三 puncture の marked Morse foliation の相対 isotopy class を指定する．
正の整数 $q_i$ と整数 $p_{ij}$，$i\ne j$，が
\begin{equation}\label{q9:four-line-arithmetic}
 \sum_{j\ne i}p_{ij}=-q_i,
 \qquad q_iq_j-p_{ij}p_{ji}=1
\end{equation}
を満たし，さらに
\begin{equation}\label{q9:four-line-basin-positive}
 B_{ij}:=w_j-w_i\frac{p_{ij}}{q_i}>0
 \qquad(i\ne j)
\end{equation}
とする．このとき，指定した isotopy classes 内で foliations の代表を選び，
本稿の比例曲率 collars と Liouville cone construction による core に，
標準 $K_4$ fan への作用適合 cone chart を構成できる．
その四本の branch denominators は $q_i$ である．
\end{proposition}

\begin{proof}
各直線の面積は $2\pi W$ であり，\eqref{q9:four-line-arithmetic} から
\[
 \sum_{j\ne i}B_{ij}=\sum_{j\ne i}w_j+w_i=W
\]
を得る．従って $L_i$ の三つの critical basins の面積を $2\pi B_{ij}$ と指定できる．
具体的には，既存の marked foliation の各 puncture disk を十分小さく固定し，
その外で正の density を調整して三 basin の面積を指定する．
全積分が変わらないので，puncture disks に相対的な Moser isotopy によって，
元の面積形式の下で所期の foliation の代表を得る．
比例曲率 collar での critical slope は
\[
 h_{ij}=\frac{w_j-B_{ij}}{w_i}=\frac{p_{ij}}{q_i}.
\]
graph-period correction は fiber tangent を消す basic form で行うので，この作用値を変えない．

$N_4=\ZZ^4/\ZZ(1,1,1,1)$ の標準 rays を $r_i$ とする．
crossing face $ij$ の regular torus の基底を，二直線の正の meridians $(\mu_i,\mu_j)$ とする．
その作用座標において
\begin{equation}\label{q9:four-line-action-matrix}
 \begin{pmatrix}I_{\mu_i}\\I_{\mu_j}\end{pmatrix}
 =\begin{pmatrix}-q_i&-p_{ji}\\-p_{ij}&-q_j\end{pmatrix}
   \begin{pmatrix}x\\y\end{pmatrix},
 \qquad F=xr_i+yr_j
\end{equation}
と定める．行列の determinant は一である．
$i$ 側の critical ray の action direction は $(-q_i,-p_{ij})$ であり，
$j$ 側は $(-p_{ji},-q_j)$ なので，両境界 rays は所期の $r_i,r_j$ へ写る．

sector 全体で $x,y>0$ であることを確認する．
$i$ 側の比例曲率 collar では $I_{\mu_i}<0$ とし，
$h=I_{\mu_j}/I_{\mu_i}$ を用いると
\[
 x=(-I_{\mu_i})(q_j-p_{ji}h),\qquad
 y=(-I_{\mu_i})(q_i h-p_{ij})
\]
である．partial basin を critical level から puncture に向かって縮めると，
$h$ は $p_{ij}/q_i$ から $w_j/w_i$ へ向かって厳密単調に増える．
従って正則部分では $y>0$ となる．
$x$ の括弧内は $h$ の affine function であり，両端でそれぞれ
\[
 q_j-p_{ji}\frac{p_{ij}}{q_i}=\frac1{q_i}>0,
 \qquad
 q_j-p_{ji}\frac{w_j}{w_i}
   =\frac{q_j}{w_i}B_{ji}>0
\]
を取るため，全 interval で正である．$j$ 側も同様である．
product crossing の極限作用点 $(-w_i,-w_j)$ における二係数は
\[
 x=q_jw_i-p_{ji}w_j=q_jB_{ji}>0,
 \qquad
 y=q_iw_j-p_{ij}w_i=q_iB_{ij}>0.
\]
transition caps は critical leaves から離れて固定できる．
そこでは上の正値性に一様な余裕があり，cap 内の normal-radius に比例する有界な補正は，
normal radius を十分小さくすることでこの余裕より小さくできる．
従って caps，product crossings，それらの丸めを通じても正の sector 内に留まる．

各 $i$ に incident な三つの maps は共通 ray 上で
$F=(-I_{\mu_i}/q_i)r_i$ に一致する．
\eqref{q9:four-line-arithmetic} の determinant 条件から
$\gcd(q_i,p_{ij})=1$ であるため，それぞれの action direction は primitive であり，
共通 branch lattice は三枚の sheets 全てで飽和している．
三 arm の period relation は，row-sum 条件 $\sum_{j\ne i}p_{ij}=-q_i$ と整合する．
標準 fan の各 cone の span lattice は $\ZZ r_i+\ZZ r_j$ なので，
\eqref{q9:four-line-action-matrix} は作用の整数余接線格子と fan の整数接線格子の
双対関係を通じて必要な格子同型を与える．

最後に，各 contact-link edge の action vector を $a(t)$ とすると，
Liouville cone の作用 map は $(r,t)\mapsto r a(t)$ である．
正則 torus fibration の作用は局所座標なので $\det(a(t),a'(t))\ne0$ であり，
正の sector 内の projective action ratio は厳密単調である\cite[Remarks 2.11--2.12]{EM}．
両 endpoint rays の間を一度だけ走るので，各 face は対応する cone sector へ同相に写る．
四つの rays 上での一致と compact link を用いて六つの faces を core まで貼れば，
定義~\ref{per:chart-definition} の作用適合 cone chart を得る．
\end{proof}

\begin{corollary}[任意の奇数分母を持つ局所 chart]\label{q9:local-all-odd}
任意の整数 $h\geq1$ に対し，$q=2h+1$ とする．
四直線を $A,B,C,D$ と標識し，
\[
 (q_A,q_B,q_C,q_D)=(q,1,1,1),\qquad
 (p_{ij})=
 \begin{pmatrix}
 0&-2h&-2&1\\
 -1&0&0&0\\
 -h&h-1&0&0\\
 2h&0&-q&0
 \end{pmatrix}
\]
とする．正の weights
\begin{equation}\label{q9:local-all-odd-weights}
 (w_A,w_B,w_C,w_D)=\left(q,h,1,1+\frac1{4h}\right)
\end{equation}
に対し，これらの分母と critical numerators を持つ局所作用適合 $K_4$ chart が存在する．
全 weights の共通の正の倍率は任意なので，任意の正の hyperplane Kähler class に対して例が得られる．
\end{corollary}

\begin{proof}
各 row の和は $-q_i$ である．$A$ に接する三辺では
$p_{ij}p_{ji}=2h=q_iq_j-1$，残る三辺では零となる．
従って全六 determinant が一である．
normalized basin areas は
\[
 (B_{ij})=
 \begin{pmatrix}
 -&3h&3&1/(4h)\\
 3h+1&-&1&1+1/(4h)\\
 3h+1&1&-&1+1/(4h)\\
 1/2&h&2h+5/2+1/(4h)&-
 \end{pmatrix}.
\]
全て正であり，各 row の和は $W=3h+3+1/(4h)$ である．
命題~\ref{q9:four-line-criterion} を適用する．
\end{proof}

この族の normal-action branch coefficients は，$A,B,C,D$ の順に
$(-3h,h,2h+2,-2)$ であり，総和は零である．
各 normal Hamiltonian や全補集合の $H_1$ が ambient linear forms に延長する条件は課していない．
三価 branch の必要条件である奇数性に対し，局所の四直線 core では任意に大きい奇数分母が生じる．
これは一つの四直線 core の幾何学的な存在族であり，五つの quintic cores の同時接着を与えるものではない．
共有 branch の分母と係数の相殺に加え，作用切片，面積，node 配分を同時に合わせる問題が残る．

% Reusable v9 lemma: can be placed after the odd-denominator K4 family.
\subsection{孤立した高分母 branch の接着障害}
\label{qh9:isolated-section}

全平面五次の一つの four-line core を考え，四 vertices の分母を
$(q,1,1,1)$，$q>1$ とする．outgoing numerators を $p_{ij}$ と書く．
各 line の三 basins が全曲線を尽くすので
\[
 \sum_{j\ne i}p_{ij}=-q_i,
 \qquad q_iq_j-p_{ij}p_{ji}=1.
\]
正の weights に対して全 basin areas が正，すなわち
$w_j-w_i p_{ij}/q_i>0$ であることも仮定する．

\begin{proposition}[唯一の高分母 vertex の係数]\label{qh9:isolated-local}
高分母 vertex から出る三 numerators は，順序を除いて
\[
 \bigl(-(q-1),-(b+1),b\bigr)
\]
である．ここで $b$ は正整数で，$b$ と $b+1$ はともに $q-1$ を割る．
$q-1=t b(b+1)$ と書くと，その vertex の normal-action 分岐係数は
\[
 \mathfrak d_{\rm high}=-t\bigl(b(b+1)+1\bigr)<0
\]
となる．
\end{proposition}

\begin{proof}
$m=q-1$ と置く．高分母 vertex を $0$ とすれば，各 edge で
$p_{0j}p_{j0}=m>0$ なので，三つの $p_{0j}$ は零でなく，絶対値は $m$ の約数である．
和は $-q=-(m+1)$ なので，二つ以上が正ではあり得ない．

全てが負で $p_{0j}=-a_j$ と仮定する．$a_1+a_2+a_3=m+1$ であり，
各 $a_j$ は正なので，いずれも $m$ には等しくない．
従って $m$ の真の約数として $a_j\le m/2$ であり，
低分母 vertex $j$ から他の二つの低分母 vertices へ出る numerators の和は
\[
 -1-p_{j0}=\frac{m}{a_j}-1>0.
\]
各低分母 vertex から少なくとも一つの正の directed edge が出る．
低分母同士では $p_{jk}p_{kj}=0$ なので，これらの正の edges は
三 vertices 上に directed cycle を持つ．正の整数 $p_{jk}$ に対して basin positivity は
$w_k>p_{jk}w_j\ge w_j$ を要求するため，cycle を一周して矛盾する．

よって正の entry がちょうど一つあり，三 entry を $(-a,-c,b)$ と書ける．
$a+c-b=m+1$ である．$a,c$ の両方が $m$ の真の約数なら
$a+c\le m$ なので不可能であり，順序を替えて $a=m$ とできる．
すると $c=b+1$ となる．$b$ と $b+1$ は互いに素でともに $m$ を割るので，
$m=t b(b+1)$，$t\in\ZZ_{>0}$ と書ける．

高分母 vertex へ入る三 numerators は
\[
 \left(-1,-\frac{m}{b+1},\frac{m}{b}\right)
\]
であり，その和は $-1+t$ である．$K_4$ fan では $\kappa=-1$ なので，
分岐係数は $-q-(-1+t)=-m-t=-t(b(b+1)+1)$ となる．
\end{proof}

\begin{corollary}[高分母 support の孤立 edge はない]\label{qh9:isolated-edge}
全平面五次の既存の branch graph $K_5$ 上に，大域作用適合 atlas があるとする．
$q>1$ の branches だけからなる部分 graph は，ただ一本の edge からなる
連結成分を持たない．
\end{corollary}

\begin{proof}
そのような edge の両端の core は，どちらも分母 pattern $(q,1,1,1)$ を持つ．
命題~\ref{qh9:isolated-local} により，その branch の両端の normal-action 係数は
ともに負となり，必要な対向端の相殺と矛盾する．
\end{proof}

これは五次の全分母が一であることまでは示さない．
二本以上の高分母 branches が隣接する場合の相殺条件と，
全共有曲線における整数 framing・実面積の同時整合が，次の問題として残る．
